\section{Calibración}

\subsection{Cálculo de errores}
\subsubsection{Definición importantes}

El valor medio de una función aleatoria continua $F(\mathbf{x})\in \mathbb{R}$ donde $\mathbf{x}:=(x_1, x_2, \ldots, x_n), x_i \in \mathbb{R}$
con una densidad de probabilidad $p(\mathbf{x})\in \mathbb{R}$,
esta definida como

\begin{equation}
  \bk{F} = \int_\Omega F(\mathbf{x}) p(\mathbf{x}) d\mathbf{x}
  \label{eq:mean}
\end{equation}

Además la suma de todas las probabilidades debe ser 1 por lo qué
\begin{equation}
  \int_\Omega p(\mathbf{x}) d\mathbf{x} = 1
  \label{eq:p-normalization}
\end{equation}

El n-esimo momento
\begin{equation}
  \bk{F^n} = \int_\Omega F^n(\mathbf{x}) p(\mathbf{x}) d\mathbf{x}
  \label{eq:moment}
\end{equation}

El n-simo momento central como
\begin{equation}
  \bk{(F^n - \bk{F})^n}
  \label{eq:n-central-momen}
\end{equation}

Y la desviación estandar como

\begin{equation}
  \sigma^2 = \bk{(F - \bk{F})^2} = \bk{F^2} - \bk{F}^2
  \label{eq:std}
\end{equation}

\subsubsection{Valor medio de un error con densidad de probabilidad uniforme}
Supongamos que se quiere calcular el error de una medición $E \in [-a, a], a > 0$,
entonces podemos definir
\begin{equation}
  E(x) = x
  \label{eq:E-error}
\end{equation}

Además que tenga igual probabilidad en cualquier punto del
interválo así $p(x) = 1/2a$ para todo el interválo, pues por la equación
\eqref{eq:p-normalization}

\begin{equation}
  \int_{-a}^a p(x) dx = 1 \implies p(x) = \frac{1}{2a}
  \label{eq:E-densidad-prob}
\end{equation}

Por lo que el valor medio de $E$ es

\begin{equation}
  \bk{E} = \int_{-a}^a x \frac{1}{2a}dx = \frac{1}{2a}\left[x\right]_{-a}^a = 0
  \label{eq:mean-error}
\end{equation}

Esto es el error medio es 0, i.e no hay sesgo si la medición es mayor o menor.


\subsubsection{Desviación estandar de un error con densidad de probabilidad uniforme}
